\section{Related Work}
\label{sec:related}

\subsection{Epigenetic clock generations and their implicit assumptions about intervention}

Epigenetic clocks are machine learning models that predict age-related quantities
from DNA methylation patterns at specific CpG sites. Three generations of clocks
have been developed, each with a different training objective and, correspondingly,
a different implicit assumption about what it means for an intervention to affect
biological aging.

First-generation clocks \cite{horvath2013,hannum2013} were trained to predict
chronological age from methylation data. They achieve high accuracy in held-out
samples ($R^2 > 0.95$) and generalise across tissues, establishing the empirical
basis for methylation-based biological age estimation. However, because they are
calibrated to chronological age, they have a structural limitation for intervention
studies: an intervention that genuinely slows biological aging would, if the clock
were perfectly measuring chronological age, show no effect. Their sensitivity to
interventions is therefore accidental rather than by design.

Second-generation clocks (PhenoAge \cite{levine2018}; GrimAge \cite{lu2019}) were
trained to predict mortality risk and health-related phenotypes rather than
chronological age directly. They show stronger and more consistent associations with
aging outcomes and are more sensitive to interventions with documented health effects.
However, they inherit a different version of the same structural problem: the training
objective does not distinguish between methylation changes that cause accelerated aging
and those that are caused by the same upstream processes that increase mortality risk.
A drug that reduces mortality by reducing inflammation, without affecting the aging
trajectory itself, would reduce a second-generation clock reading by design.

Third-generation measures, notably DunedinPACE \cite{belsky2022} and DunedinPoAm38,
were trained on longitudinal biomarker data to estimate the \emph{pace} of aging
rather than biological age as a static quantity. These measures show the strongest
and most consistent sensitivity to interventions in clinical trials, including caloric
restriction \cite{belsky2023calerie}. We argue in Section~\ref{sec:framework} that
this superior sensitivity reflects a structural fact: pace-of-aging measures estimate
the rate of traversal of the aging trajectory ($\dot{s}$), while static clocks
estimate position ($s$). These are genuinely different quantities that an intervention
can affect independently. The apparent superiority of DunedinPACE for intervention
studies is not simply better calibration, but reflects what these measures are
actually measuring.

Across all three generations, a common assumption is implicit but never stated: that
a clock's output is a valid measure of biological aging by virtue of its training
objective. This makes intervention evaluation dependent on clock choice, which is the
circularity we aim to resolve.

\subsection{The reliability and false-positives literature}

A recent and practically important line of work has addressed the problem of false
positive intervention effects in epigenetic clock studies. Borrus et al.\
\cite{borrus2024} systematically evaluated clock responses to a range of interventions
and showed that first-generation clocks frequently produce effects that are not
replicable across high-reliability clock variants, do not survive multiple testing
correction, and are inconsistent with effects of second- and third-generation clocks
on the same datasets. They conclude that mortality-trained clocks provide more reliable
intervention signals.

This work identifies a real problem, and its empirical findings are directly relevant
to our framework. However, its criterion for a genuine intervention effect remains
statistical: an effect is genuine if replicable across multiple clocks and surviving
multiple testing correction. This criterion is necessary but not sufficient. A
replicable, multi-clock-consistent intervention effect can still be clock-gaming in
our sense: if an intervention systematically displaces methylation in a direction
orthogonal to the aging trajectory, multiple clocks whose weight vectors have
components in that direction will consistently report an effect, even though the
cell's position on the aging trajectory has not changed.

The distinction is not merely theoretical. Any compound specifically designed or
selected to reduce methylation at clock CpG sites would produce replicable,
multi-clock-consistent effects by construction. The statistical reliability criterion
cannot detect this; the geometric criterion we propose can.

TranslAGE \cite{translagee2025}, a recently published platform harmonising 179 human
blood methylation datasets with pre-calculated scores for 41 epigenetic clocks,
provides the infrastructure for large-scale empirical evaluation. We use its
harmonised datasets and accession list as the primary data source for the empirical
component of this paper.

\subsection{Causal enrichment approaches}

The most technically sophisticated prior work on grounding epigenetic clocks in
causal structure is the causality-enriched epigenetic age framework of Ying et al.\
\cite{ying2024}. Using epigenome-wide Mendelian randomization, they identified CpG
sites with potential causal links to aging-related traits, then trained separate
clocks on damage-associated sites (DamAge) and adaptation-associated sites (AdaptAge).
DamAge correlates with adverse outcomes including mortality; AdaptAge correlates with
beneficial adaptations.

We treat this work as complementary rather than competing. Their key contribution is
distinguishing causally upstream methylation changes from downstream correlates.
However, the framework remains clock-centric in a specific sense: it selects better
CpG features and trains clocks on them, rather than identifying the aging manifold
independently of any clock. The causal enrichment criterion is a property of
individual CpG sites, not of the geometric structure of methylation space as a whole.

Our framework makes a testable prediction about the relationship between these
approaches: the causally enriched CpG sites identified by Ying et al.\ should have
weight vectors more closely aligned with our aging direction $v^*$ than standard
clock CpG sites. We test this prediction in Section~\ref{sec:results}.

\subsection{Statistical assumptions of biological age predictors}

Sluiskes et al.\ \cite{sluiskes2024} provide a careful analysis of the statistical
assumptions implicit in widely used biological age predictors, including multiple
regression, the Klemera-Doubal method, and PCA-based approaches. They show that these
assumptions are frequently violated and that violations have quantifiable consequences
for biological age estimates.

This work is complementary to ours on the theoretical side. Where Sluiskes et al.\
examine what existing methods statistically assume, we examine what a valid definition
of biological aging should formally require. The desiderata we derive in
Section~\ref{sec:framework} can be seen as the normative counterpart to their
descriptive analysis of assumption violations.

\subsection{Multi-omic clocks and the substrate question}

A growing literature has developed biological age measures from proteomics
\cite{lehallier2019}, metabolomics, glycan profiles, and combinations thereof.
Studies of individual aging trajectories have shown that individuals differ in their
ageotype --- which molecular substrate ages fastest in a given person \cite{shen2020}
--- suggesting aging is not a single process tracked equally by all substrates.

We develop our framework within the methylation substrate specifically. Methylation
has structural properties that make it appropriate for the cross-context identification
criterion we propose: it is measured on a common platform across tissues, enabling
cross-tissue displacement comparisons at the CpG level; it is positioned causally
upstream of many other aging hallmarks; and it is responsive to interventions on
timescales accessible to clinical trials. Other molecular clocks are used in our
validation as independent evidence for the biological relevance of our trajectory
estimate, not as inputs to its construction. We note explicitly that concordance
between imperfect proxies does not validate any individual proxy; mortality prediction
on held-out data is the appropriate validation criterion.

\subsection{The circularity problem}

Moqri et al.\ \cite{moqri2023} provide the most explicit acknowledgment in the
literature of the circularity inherent in using epigenetic clocks to evaluate
interventions. Their Figure~4 diagrams the relationship: biomarkers of aging are used
to identify geroprotectors, while geroprotectors are used to validate biomarkers.
They propose using a small set of interventions with strong prior evidence of
geroprotective effects to benchmark clocks, effectively using interventions to
validate clocks rather than vice versa.

This is a practically useful step, but it relocates rather than resolves the
circularity: the prior-validated geroprotectors must themselves be grounded in
something other than clock readings, and what that grounding is remains informal.
Our framework makes the grounding explicit. The labeling sets $\mathcal{I}^+$ and
$\mathcal{I}^-$ are grounded in mortality and healthspan evidence independent of
methylation clocks --- epidemiological evidence for smoking, randomised trial evidence
for caloric restriction, multi-species lifespan evidence for dietary restriction.
The residual circularity in our framework --- that biological age is ultimately
anchored to health outcomes --- is unavoidable and appropriate: any definition of
biological age disconnected from health outcomes would be biologically uninteresting.
What we eliminate is the clock-dependence.
